% Completion of Kamp's Theorem III.2 from the full ten-lemma interval construction.

\subsubsection{Completion of the first-order-to-temporal construction}
\label{sec:p09-kamp-iii2-completion}

The final step converts the regional components in the one-variable
standard form into elementary one-variable formulas, then follows their
elementary construction histories to temporal formulas. The initial
reduction on source pp.~60--62 and the conclusion on pp.~93--94 are
both needed: a general Bef or Aft component can contain nested regional
components and can begin with either quantifier
\cite[ch.~III, pp.~60--62, 93--94]{Kamp1968}.

\begin{hblemma}{thm:p09-kamp-one-sided-closure}{proved=annex:p9.kamp.one.sided.closure}{Closing a supplied interval representation on one side}
Let \(x,w\) be distinct variable names.
If \(\rho(x,w)\) has a supplied \(I(x,w)\) representation,
then \(\exists x\,(x<w\land\rho(x,w))\) has an actual
\(\operatorname{E}(w)\) presentation equivalent over
\(\mathsf{Com}\). If \(\rho(w,x)\) has a supplied
\(I(w,x)\) representation, its guarded future existential
\(\exists x\,(w<x\land\rho(w,x))\) has such a presentation.

For the past case, distribute the guarded existential over the given
elementary disjuncts and use the elementary past-closure construction.
For the future case first convert each elementary past word into its
same-word elementary future presentation, then close it on the future
side. Retain every disjunct. The empty disjunction gives \(F^K_w\).
\end{hblemma}

\begin{hbdefinition}{def:p09-kamp-primitive-regions}{The source's primitive regional class}
Let \(\mathcal D(w)\) consist of the members of
\(\operatorname{Bef}(w)\cup\operatorname{Aft}(w)\) that have
no proper subformula occurrence belonging to
\(\operatorname{Bef}\cup\operatorname{Aft}\). This is the
source's class \(\mathcal D\), with the sole free variable
displayed. Its members begin with an existential or universal regional
quantifier. Their matrices contain only At conditions at the quantified
point and Bet conditions between it and the fixed endpoint.

An \emph{existential envelope} also permits finitely many At conditions
at the fixed endpoint. Its two forms are
\[
 \exists x\,(x<w\land\Theta(x,w)),\qquad
 \exists x\,(w<x\land\Theta(w,x)),
\]
where \(\Theta\) is a finite conjunction of Bet components and
At components at either endpoint. This supplementary envelope notation
records the source's final display~(28); it does not enlarge the
syntactic definition of \(\mathcal D\).
\end{hbdefinition}

\begin{hblemma}{thm:p09-kamp-primitive-regions}{proved=annex:p9.kamp.primitive.regions}{Elementary presentations of all primitive regional cases}
Every existential envelope and every member of \(\mathcal D(w)\)
has an effectively constructed \(\operatorname{E}(w)\)
presentation equivalent over \(\mathsf{Com}\).

For an existential envelope, apply Lemma~10 to its matrix, then use
the appropriate one-sided closure. For a universal primitive member,
negate the universal, normalize its complemented interval matrix by
\cref{thm:p09-kamp-interval-boolean-normalization}, keep all resulting
finite disjunctive alternatives, and close their existential envelopes.
Negate the resulting elementary formula to recover the original universal.
The construction does not treat a finite disjunction of primitive
formulas as one primitive formula by definition
\cite[ch.~III, pp.~61--62, 93--94]{Kamp1968}.
\end{hblemma}

\begin{hbmethod}{def:p09-kamp-regional-restoration}{Restoring maximal proper regional occurrences}
For a Bef or Aft formula \(\varphi(w)\), select all maximal proper
subformula occurrences \(\varphi_1(y_1),\ldots,\varphi_k(y_k)\)
that belong to \(\operatorname{Bef}\cup\operatorname{Aft}\).
Maximal means not contained in any larger proper occurrence of that kind.
These are occurrences in the finite syntax tree, including distinct
occurrences with identical text. They are disjoint.

Recursively obtain an elementary presentation for each selected occurrence.
Choose distinct fresh unary letters \(Q_1,\ldots,Q_k\), absent
from the original formula and every obtained presentation, and replace
the selected occurrence \(\varphi_i(y_i)\) by \(Q_i(y_i)\).
The resulting skeleton belongs to \(\mathcal D(w)\).
Construct its elementary presentation using the preceding theorem.
Properly substitute the obtained elementary formula for each \(Q_i\),
renaming its free argument and bound variables as required.
All target letters are absent from all replacements. If \(k=0\),
the formula already belongs to \(\mathcal D(w)\).
\end{hbmethod}

\begin{hblemma}{thm:p09-kamp-regional-elementary}{proved=annex:p9.kamp.regional.elementary}{Elementary presentations of all past and future regional components}
Every member of \(\operatorname{Bef}(w)\cup\operatorname{Aft}(w)\)
has an effectively constructed \(\operatorname{E}(w)\)
presentation equivalent over \(\mathsf{Com}\). The construction
just given terminates by recursion on proper original subformula occurrences.
It includes existential and universal roots, repeated occurrences, nested
quantified arguments, and both temporal directions. All substitutions
preserve the intended binding of the sole free argument of each replaced
regional component \cite[ch.~III, pp.~60--62, 93--94]{Kamp1968}.
\end{hblemma}

\begin{hbdefinition}{def:p09-kamp-translation-range}{The image of the temporal translation}
Keep the source temporal language \(\mathsf{TL}_1\) and its
first-order translation \(D\), with distinguished current-point
variable \(t_0\). Define
\[
             \mathcal C\defeq\{D(\eta):\eta\in\mathsf{TL}_1\}.
\]
This is the source's image class, distinct from
\(\operatorname{E}\). An intermediate elementary formula need
not literally belong to \(\mathcal C\); its constructed temporal
expression \(\eta\) supplies the actual image member \(D(\eta)\).
For another free variable \(w\), use capture-avoiding replacement
of the free \(t_0\) by \(w\), after freshening bound names.
\end{hbdefinition}

\begin{hbmethod}{def:p09-kamp-final-translation}{The complete first-order-to-temporal procedure}
Input an \(L\)-formula \(\varphi\) with exactly one free variable
\(w\). Construct its one-variable standard form by Theorem~III.1.
Its sole-variable prefix is true over linear orders; replace that
prefix by the source truth label. Its other components belong to
\(\operatorname{At}(w)\), \(\operatorname{Bef}(w)\) or
\(\operatorname{Aft}(w)\).

Keep each At component as an elementary seed. Replace each Bef and Aft
component by the elementary presentation just constructed, and keep
the standard form's finite conjunctions and disjunctions. This gives
an actual \(\operatorname{E}(w)\) construction history.
Apply the elementary-to-temporal procedure of Lemma~3 to that history.
Return its \(\mathsf{TL}_1\) formula \(\eta\) and the literal
image member \(D(\eta)\). Every proper-substitution placeholder
has been eliminated before this final output.
\end{hbmethod}

\begin{hbtheorem}{thm:p09-kamp-iii2}{proved=annex:p9.kamp.iii2}{Kamp's expressive-completeness theorem over complete linear orders}
For every \(L\)-formula \(\varphi(w)\) with exactly one free
variable, the stated procedure constructs a temporal formula
\(\eta\in\mathsf{TL}_1\) such that
\[
          \mathsf{Com}\models
                \varphi(w)\longleftrightarrow D(\eta)[w/t_0].
\]
In particular every tense-defining formula at \(t_0\) is
\(\mathsf{Com}\)-equivalent to a member of \(\mathcal C\),
the source's Theorem~III.2. One formula works uniformly for every
complete linear order, every unary interpretation and every current
point. Density and absence of endpoints are not assumed.

The temporal constructors are exactly negation, conjunction and the
source's strict \(S,U\), with their witness operand first and
their open-interval operand second. Thus they express every uniformly
first-order-definable tense over \(\mathsf{Com}\).
For a positive finite arity \(n\), if the allowed unary signature
is \(P_0,\ldots,P_{n-1}\), the output can use only the corresponding
\(q_0,\ldots,q_{n-1}\). The literal source constant convention
may use \(P_0\) even when it was absent from the input formula.
The proof establishes an effective terminating construction; no
output-size or running-time bound is asserted here
\cite[ch.~III, thm.~2, p.~60 and pp.~93--94]{Kamp1968}.
\end{hbtheorem}

\begin{hbexample}{def:p09-kamp-final-universal-example}{A universal past condition and its finite disjunction}
Let \(Q,R\) be unary letters, with temporal counterparts \(q,r\).
For
\[
 \varphi(w)=\forall z\,(z<w\to\neg Q(z)\land\neg R(z)),
\]
the complemented universal gives
\[
 \neg\varphi(w)\longleftrightarrow
       \exists z\,(z<w\land Q(z))\ \lor\
       \exists z\,(z<w\land R(z)).
\]
With \(\tau=\neg(q\land\neg q)\), a temporal expression
for \(\varphi\) is
\[
                        \neg S(q,\tau)\land\neg S(r,\tau).
\]
At a first point both strict past existentials are false, so this
formula is true, matching universal vacuity. If the sole predecessor
in a two-point order satisfies either \(Q\) or \(R\), it is
false. The disjunctive alternatives produced during normalization
must therefore be retained.
\end{hbexample}

\paragraph{End of the source construction.}
The concluding argument ends on printed p.~94. Chapter~IV begins on
p.~95 with other expressibility results. The ten interval lemmas and
the preliminary reduction now supply the complete constructive argument
for the stated Theorem~III.2.
